\section{Policy Implications}
\label{sec:policy}

The model delivers concrete, actionable implications for three areas of
economic policy.

\subsection{Monetary Policy}

\subsubsection{Augment the Taylor Rule}

The central result---Proposition 3---prescribes augmenting the standard
Taylor rule with two additional terms:

\begin{enumerate}
  \item \textbf{Labor market concentration indicator} ($\phi_\mu^* > 0$):
    Central banks should ease more aggressively when real-time measures
    of labor market concentration indicate deepening monopsony. Operationally,
    this could be implemented using quarterly HHI data from online job
    postings \citep{azar2020} or administrative employment records.
    Our calibration suggests $\phi_\mu^* \approx 0.35$ in the welfare-optimal
    rule---a modest but quantitatively significant adjustment.

  \item \textbf{Supply chain fragmentation index} ($\phi_\phi^* < 0$):
    Central banks should tighten less aggressively when fragmentation shocks
    drive inflation, as these are supply-side phenomena that rate hikes cannot
    efficiently address. Operationally, this could condition on indices of
    trade policy uncertainty, cross-bloc trade flow divergence, or the
    shipping cost indices that proved to be leading indicators of supply
    disruptions in 2021.
\end{enumerate}

\subsubsection{Revise the Natural Rate Concept}

The model implies that the ``natural rate of unemployment''---the rate
consistent with stable inflation---is not structural but depends on current
market structure. Under monopsony, the unemployment rate consistent with
wage stability (the wage-consistent NAIRU) is lower than the rate consistent
with price stability (the inflation-consistent NAIRU). These two concepts
diverge by:
\begin{equation}
  u^*_{price} - u^*_{wage} = \frac{\bar\mu^w}{\gamma}\left[\frac{1}{\varepsilon_p} -
  \frac{1}{\varepsilon_w}\right]
  \label{eq:nairu_gap}
\end{equation}
In our calibration, this gap is approximately 0.4 percentage points. Central
banks using wage-consistent NAIRU estimates (as in the \citet{phillips2021}
wage-price spiral literature) may systematically overtighten relative to
price stability, at the cost of unnecessary unemployment.

\subsubsection{Forward Guidance Under Monopsony}

The model implies that the sacrifice ratio is state-dependent and higher
in recessions (when monopsony deepens). This has implications for forward
guidance: promises of future accommodation should be calibrated more
aggressively in recessions than standard models suggest, because the
output cost of premature tightening is larger. Conversely, the benefit of
credible inflation anchoring is amplified: if households believe inflation
will revert, real wages are not compressed by the markdown channel, partially
offsetting the welfare cost of the monetary contraction.

\subsection{Antitrust and Labor Market Policy}

A novel finding of the model is that \emph{antitrust policy is
macroeconomic policy}. Reducing labor market concentration by enforcing
competitive labor markets raises $\bar\varepsilon$, which:
\begin{enumerate}[(i)]
  \item Steepens the Phillips curve (Proposition 1), making monetary policy
    more effective;
  \item Reduces the countercyclical variation in $\mu_t^w$, stabilizing
    real wages and aggregate demand over the cycle;
  \item Raises the labor share and aggregate welfare directly.
\end{enumerate}
A policy that raises $\bar\varepsilon$ from 4 to 6 (the ``mild monopsony''
case) would steepen the Phillips curve slope from 0.029 to 0.038 (+31\%) and
reduce the sacrifice ratio from 34.4 to 26.3 ($-24\%$). The associated
reduction in the output cost of the 2022--2023 disinflation would have been
approximately $0.24 \times 2\% = 0.48$ percentage points of output---worth
roughly \$120 billion in 2023 GDP terms.

\subsection{Industrial and Trade Policy}

\subsubsection{Managed Reshoring vs.\ Abrupt Decoupling}

The fragmentation TFP loss formula \eqref{eq:tfp_loss} provides a precise
cost function for trade policy design. Key implications:

\begin{enumerate}
  \item \textbf{The speed of decoupling matters.} A gradual increase in
    $\phi_t$ allows firms to build domestic substitutes before cross-bloc
    links are severed, reducing peak Domar weight reallocation. Our model
    quantifies the annual TFP cost of abrupt decoupling at 1.2\% vs.\ 0.4\%
    for a ten-year managed transition---a 3:1 cost ratio.

  \item \textbf{Sector-specific intervention is more efficient than uniform tariffs.}
    Sectors with high upstream Domar weights and limited domestic substitution
    capacity (Energy, Semiconductors, Critical Minerals) are the
    highest-priority targets for supply-chain resilience investment. Uniform
    tariffs that raise $\phi_t$ across all sectors impose disproportionate
    costs through these concentrated upstream nodes.

  \item \textbf{Strategic buffer stocks are welfare-improving.}
    In a fragmented world, maintaining inventories of inputs with high cross-bloc
    Domar weights effectively reduces the volatility of $\phi_t$ from the
    perspective of domestic producers. Our model provides a formula for the
    optimal inventory level: it should equal the expected peak TFP loss during
    a fragmentation episode divided by the cost of carry. This provides a
    formal justification for the CHIPS and Science Act's semiconductor
    stockpile provisions and the IRA's critical minerals investments.
\end{enumerate}

\subsubsection{Complementarity Between Monetary and Industrial Policy}

A broader insight is that monetary and industrial policy are
\emph{complements}, not substitutes, in the fragmented-monopsonistic world
we model. Monetary policy can smooth the business cycle; it cannot change
the underlying structure of labor markets or supply chains. Industrial
policy can reduce $\bar\phi$ (by building domestic capacity) and raise
$\bar\varepsilon$ (by fostering competition); it operates at the wrong
frequency to smooth short-run fluctuations. The welfare optimum requires both
instruments working in tandem, with the central bank conditioning on the
current state of structural policy through the augmented Taylor rule.
